\documentclass[11pt,a4paper]{article}
\usepackage[margin=1in]{geometry}
\usepackage{hyperref}
\usepackage{enumitem}
\usepackage{xcolor}
\usepackage{listings}

\lstset{
  basicstyle=\ttfamily\small,
  breaklines=true,
  frame=single,
  columns=fullflexible,
}

\title{\bfseries Week 3: Real GitHub OAuth Evidence}
\author{Saksham Boran \\ \small Roll No: 1024030856}
\date{12 September 2026}

\begin{document}
\maketitle

\section{Context}
Week 2's profile page displayed whatever GitHub URL a user had
typed in during onboarding -- an unverified string, not evidence.
This week replaced that placeholder with a real GitHub OAuth
connection, reused across the dashboard, onboarding's GitHub step,
and the profile page, whose result merges directly into the user's
skill ratings. The work spanned a hand-rolled OAuth flow, a
database transaction, and a reconciliation rule shared with the
existing resume-upload evidence path.

\section{Backend: Keeping the Schema in Sync}
Between this and the prior weeks, \texttt{Profile.skillRatings}
and \texttt{Profile.activeResumeId} had been added to the schema
by migrations landing alongside the resume-parsing work. Before
starting on GitHub evidence, pulled those migrations and
re-synced my local Neon database (\texttt{prisma migrate deploy} /
\texttt{prisma generate}) so the Prisma client and the actual
database agreed on the shape of \texttt{Profile} before writing
any code against it.

\section{Backend: Hand-Rolled GitHub OAuth}
GitHub was never a sign-in provider -- Week 2's Google-only,
college-domain-restricted sign-in stayed as the only auth method.
GitHub is purely an \emph{evidence} source, connected after a user
is already signed in, so it stays outside NextAuth entirely and is
implemented as two plain route handlers.

\subsection{Design: Two API Routes}
\begin{itemize}[leftmargin=0.7cm]
  \item \textbf{\texttt{GET /api/github/start}} -- generates a
    random CSRF \texttt{state} value, stores it in an
    \texttt{httpOnly} cookie alongside a \texttt{returnTo} value
    (also cookie-stored), and redirects to GitHub's OAuth
    authorize URL.
  \item \textbf{\texttt{GET /api/github/callback}} -- GitHub
    redirects here with \texttt{?code\&state}. The route validates
    \texttt{state} against the cookie, exchanges \texttt{code} for
    an access token server-side (the client secret never touches
    the browser), fetches the user's public repos, derives skills
    from the repos' primary languages, and writes the result to
    the database in one transaction.
\end{itemize}

\subsection{Security Detail: The \texttt{returnTo} Allowlist}
``Connect GitHub'' appears in three places (dashboard, onboarding,
profile), and after the OAuth round-trip the callback needs to
send the user back to whichever one they started from. Rather than
trusting an arbitrary \texttt{returnTo} query value -- an open
redirect if done naively -- it's checked against a fixed allowlist
before being used to build the final redirect:
\begin{lstlisting}[language=JavaScript]
const RETURN_TO_ALLOWLIST = ["/dashboard", "/profile", "/onboarding"];

export function sanitizeReturnTo(value, fallback) {
  if (value && RETURN_TO_ALLOWLIST.includes(value)) {
    return value;
  }
  return fallback;
}
\end{lstlisting}

\subsection{Deriving Skills from Public Repos}
\texttt{fetchPublicRepoLanguages} calls the GitHub repos API for a
connected user's own (non-fork) repositories and reads each repo's
single primary \texttt{language} field -- a full byte-count
breakdown per repo was considered and rejected, since it would
need one extra API call per repository for marginal accuracy.
Forks are excluded so a starred/forked repo doesn't misattribute
someone else's work to the connecting user.

\texttt{deriveSkillsFromLanguages} then ranks languages by
frequency and maps GitHub's linguist language names onto the
platform's fixed skill-slug list
(\texttt{GITHUB\_LANGUAGE\_TO\_SKILL} in \texttt{lib/skills.ts}).
Languages with no corresponding skill slug (e.g.\ ``HTML'',
``Shell'') are deliberately dropped rather than guessed at, and
only true linguist-detected languages are mapped here --
frameworks and tools like React or Docker aren't detectable this
way and are left to the resume-parsing path instead.

\section{Backend: Reconciling Skill Ratings Across Two Evidence
  Sources}
The core design decision this week: skill ratings should be
\emph{reconciled}, not overwritten, regardless of which evidence
source produces them. \texttt{mergeSkillRatings} (moved from
\texttt{lib/resume-upload.ts} into the shared
\texttt{lib/skills.ts}, since it's now used by both the resume and
GitHub paths) only fills in skills the user doesn't already have a
rating for:
\begin{lstlisting}[language=JavaScript]
export function mergeSkillRatings(existing, newSkills) {
  const merged = { ...existing };
  for (const slug of newSkills) {
    if (!merged[slug]) {
      merged[slug] = "INTERMEDIATE";
    }
  }
  return merged;
}
\end{lstlisting}
This means a resume-derived rating and a GitHub-derived rating for
the same skill never fight each other, and a user's own
self-reported rating (if they already had one from Path B
onboarding) is never silently downgraded by a new evidence source.

\subsection{Transaction: One Evidence Record per Account}
Unlike resumes, where a user can have several historical uploads,
a GitHub connection is treated as singular -- reconnecting
replaces rather than accumulates. The callback route wraps the
delete-then-create-then-upsert sequence in a single
\texttt{prisma.\$transaction}:
\begin{lstlisting}[language=JavaScript]
await tx.evidenceRecord.deleteMany({ where: { userId, source: "GITHUB" } });
const created = await tx.evidenceRecord.create({
  data: { userId, source: "GITHUB", payload },
});
await tx.profile.upsert({
  where: { userId },
  create: { userId, githubUsername, skillRatings: mergedSkillRatings },
  update: { githubUsername, skillRatings: mergedSkillRatings },
});
\end{lstlisting}
This also cleanly clears out the old unverified-URL placeholder
record the first time a user connects for real, and the
transaction guarantees the evidence record and the merged skill
ratings never end up out of sync if one write fails partway
through.

\section{Frontend Consequence: Surviving a Full-Page Redirect}
Because OAuth requires a full-page navigation away to GitHub and
back, any in-progress client state is lost by default. Onboarding's
multi-step wizard needed its state persisted to
\texttt{sessionStorage} before redirecting, so a user connecting
GitHub mid-onboarding returns to the same step with their other
answers intact rather than starting over. Also built a shared
\texttt{ConnectGithubButton} used identically across the dashboard,
onboarding's GitHub step, and the profile page, rather than three
slightly different implementations of the same start-the-OAuth
click handler.

\section{Verification}
No automated test framework exists yet, so verification was
manual: TypeScript compiles cleanly, the production build
succeeds, and a live walkthrough against the real Neon database
and a real GitHub OAuth App confirmed:
\begin{itemize}[leftmargin=0.7cm]
  \item Connecting GitHub from the profile page, dashboard, and
    onboarding's GitHub step all round-trip correctly and land
    back on the originating page.
  \item A connection creates exactly one \texttt{GITHUB}
    \texttt{EvidenceRecord} and merges derived language skills into
    \texttt{Profile.skillRatings} without touching any rating the
    user or a resume upload had already set.
  \item Reconnecting a different GitHub account replaces the old
    evidence record rather than leaving two.
\end{itemize}

\section{Outcome}
Shipped the full GitHub OAuth evidence connection -- covering
\texttt{lib/github.ts}, both API routes, the shared
\texttt{ConnectGithubButton}, and the skill-ratings reconciliation
now shared with the resume path -- replacing the onboarding-only
unverified-URL placeholder everywhere it appeared.

\section{Follow-ups}
\begin{itemize}[leftmargin=0.7cm]
  \item GitHub-derived skills default to \texttt{INTERMEDIATE}
    like resume-derived ones; there's no signal yet (e.g.\ repo
    count or recency) that would justify a different starting
    level.
  \item The public-repos scan is capped at 100 repos per user
    (\texttt{MAX\_REPOS\_SCANNED}); untested against an account
    with more than that.
  \item Access tokens are used once during the callback and
    discarded -- nothing is stored long-term, so re-deriving skills
    later (e.g.\ on a schedule) would require the user to
    reconnect.
\end{itemize}

\end{document}
